\documentclass[a5paper,10pt]{memoir}
% Packages
\usepackage[nobottomtitles]{titlesec}
\usepackage{microtype}
\usepackage[yyyymmdd,hhmmss]{datetime}
%\usepackage[british, czech]{babel}
\usepackage{polyglossia}
\setdefaultlanguage{english}
\setotherlanguages{czech, hebrew}
\usepackage{xcolor}

%!TEX TS-program = lualatex 
%!TEX encoding = UTF-8 Unicode
%\usepackage{fontspec}
%\setmainfont{TeX Gyre Pagella} %palatino clone
% local
\input{rubrics.sty}
\input{styling.sty}
\input{titlepage.sty}
% Config
\copypagestyle{versionplain}{plain}
\newcommand{\version}{\tiny Release: v1.0.0 (template v2.3.0)}
\makeevenfoot{versionplain}{\thepage}{}{\version}
\makeoddfoot{versionplain}{\version}{}{\thepage}
\pagestyle{versionplain}
\thispagestyle{empty}
% =========
% per-document styling: we don't set this in `styling.sty' as we might want
% different layouts.
% =========
% layout
\setulmarginsandblock{0.5in}{0.7in}{*}
\setlrmarginsandblock{0.5in}{0.5in}{*}
\checkandfixthelayout
% headers
\setsecnumdepth{chapter}
\setsecheadstyle{\Large\scshape\raggedright\centering}
\setsubsubsecheadstyle{\color{darkblue}\large\scshape\raggedright\centering}
\setparaheadstyle{\color{red}\large\raggedright\centering}
\setafterparaskip{1.5ex plus .2ex}
\setlength{\parskip}{\medskipamount}
\setlength{\parindent}{0em}
\midsloppy

% REDEFINE CHAPTERS AND SECTIONS
\usepackage{xcolor}
% define colors
\definecolor{sectioncolour}{RGB}{23,85,142}
\definecolor{chaptercolour}{RGB}{0,51,102}
\definecolor{subsectioncolour}{RGB}{23,85,142}

\setsecheadstyle{\Large\bfseries\color{sectioncolour}}% Set \section style
\setsubsecheadstyle{\large\bfseries\color{subsectioncolour}}% Set \subsection style

\renewcommand{\chapnamefont}{\normalfont\large\scshape\raggedleft\color{chaptercolour}}
\renewcommand{\chapnumfont}{\Huge\bfseries\color{chaptercolour}}
\renewcommand{\chaptitlefont}{\normalfont\Huge\bfseries\sffamily\raggedleft\color{chaptercolour}}

\usepackage{multicol}

\input{couching}
% This file defines the blessings to be inserted into haggadah.tex  In this way
% multiple formats can be typeset very quickly.  haggadah.tex should not normally
% need changing.  Note that this is not the most readable way to insert text
% into a LaTeX document, but it is the most powerful: the macros defined here
% are directly excecuted when building the document.

%!TEX TS-program = lualatex 
%!TEX encoding = UTF-8 Unicode
%\usepackage{fontspec}
\usepackage{polyglossia}
\setdefaultlanguage{english}
\setotherlanguages{czech, hebrew}
%\newfontfamily{\hebrewfont}{New Peninim MT}
% 

\newcommand{\CandlelightingBlessingTransliterated}{%
	{Baruch atah Adonai eloheinu melech ha-olam asher kidd’shanu
	b’mitzvotav v’tzivvanu l’hadlik ner shel (shabbat v’shel) yom tov.}
\newcommand{\CandlelightingBlessingEnglish}
	{ May you be blessed, Lord our G-d, ruler of the universe, who sanctifies us via his commandments and who commands us to light the lights of Shabbat and festivals.}
}

\newcommand{\WineBlessingHebrew}{%
	\texthebrew{שלום עולם
		ברָּ וך ְּ אַ תָּ ּה יי ,ָּ אֱֹלהֵ ינו מֶּ לךֶּ ְּ הָּ עֹולם,ְָּּ בוֹּ רֵ א פ רִ י הַ גפֶָּּ ן
		}
	}
	
\newcommand{\WineBlessingTransliterated}{%
		{Baruch atah Adonai eloheinu melech ha-olam borei p’ri ha-gafen.}}
		
\newcommand{\WineBlessingEnglish}
		{We thank You Lord our God, ruler of the universe, for making the fruit of the vine.}

\newcommand{\ShehechyanuBlessing}{%

We praise You, Eternal God, Ruler of the universe, for keeping us alive,
sustaining us, and enabling us to reach this season.
}

% 
\newcommand{\KarpasBlessingTransliterated}{%
	Baruch atah Adonai eloheinu melech ha-olam borei p’ri ha-adamah.
}
\newcommand{\KarpasBlessingEnglish}{%
	We thank You God, for making food grow from the ground.
}
% 
\newcommand{\HaMotziBlessingTransliterated}{%
	Baruch ata Adonai Eloheinu melech ha’olam hamotzi lechem min ha’aretz.}
\newcommand{\HaMotziBlessingEnglish}{Blessed are you, Lord our God, ruler of the universe who brings forth bread from the earth.}
		
\newcommand{\HaLachmaTransliterated}{%
		
		Ha lachma anya di achalu
		av’hatana b’ar’ah d’mitzrayim.
		Kol dich’fin yeiy’tei v’yeichul
		Kol ditz’rich yeiy’tei v’yif’sach.
		Ha-shata hacha –
		l’shata d’atya b’ar’ah d’yisra’el
		Ha-shata avdei –
		l’shata d’atya b’nei chorin}
	
\newcommand{\HaLachmaEnglish}{%
		This is the bread of affliction our ancestors ate in the land of Egypt. Let all who are hungry come and eat; let all who are in need come and share our Passover. This year here, next year in the land of Israel; this year oppressed, next year free.}
	
\newcommand{\MaNishtanaText}{%
	Mah nishtanah ha-lailah ha-zeh mi-kol ha-leilot!
	She-b’chol ha-leilot anu ochlin chameitz u-matzah, ha-lailah ha-zeh kulo matzah?
	She-b’chol ha-leilot anu ochlin sh’ar y’rakot, ha-lailah ha-zeh maror?
	She-b’chol ha-leilot ein anu matbilin afilu pa’am achat, ha-lailah ha-zeh sh’teif’amim?
	She-b’chol ha-leilot anu ochlin bein yoshvin u-vein m’subin, ha-lailah ha-zeh
	kullanu m’subin?
}
\newcommand{\MaNishtanaTextTranslation}{%
		How different this night is from all other nights!
	On all other nights we eat either leavened or unleavened bread…why only unleavened bread tonight?
	On all other nights we eat different types of herbs and vegetables…why bitter herbs tonight?
	On all other nights we do not even dip once…why do we dip twice tonight?
	On all other nights we eat either sitting or leaning…why do we all lean tonight?}
	
\newcommand{\MaggidText}{%
	\l{%
		passage or poem for story
	}
	\e{%
		We present to thee, O Lord, these gifts for our sacrifice, as honoring thy holy ones may they please thee, and of thy pity, may they become profitable to us.
	}
}

\newcommand{\PlaguesText}{%
	\begin{itemize}
		\color{sectioncolour}
		\item Dam
		\item Tz'farde'a
		\item Kinim
		\item A'rov
		\item Dever
		\item Sh'chin
		\item Barad
		\item Ar'beh
		\item Choshech
		\item Makkat b'chorot
	\end{itemize}}
	
\newcommand{\DayyeinuResponsa}{%
	\begin{itemize}
		If G-d had brought us out of Egypt and not given us Shabbat : Dayyeinu
		\end{itemize}
	}
%
\newcommand{\HallelBlessing}{%
	\l{%
		Hallelu hallelu hallelu, hallelu, halleluyah!
		Kol ha-n’shamah t’hallel yah, hallelu halleluyah!
		Let us give praise – Let us all praise God. Halleluyah!
	}
	
\newcommand{\NirtzahBlessing}{%
	\l{%
		Next year in jerusalem findtext
	}

	
	\e{%
	\per
}





% File paths: we don't use symlinks as (a) not all platforms support them, and
% (b) they don't fit nicely with the flow we're using.

\newcommand{\kyriePath}{../Ordinaries/masses/4/kyrie}
\newcommand{\gloriaPath}{../Ordinaries/masses/4/gloria}
\newcommand{\sanctusPath}{../Ordinaries/masses/4/sanctus}
\newcommand{\agnusPath}{../Ordinaries/masses/4/agnus}
\newcommand{\itePath}{../Ordinaries/masses/4/ite}
% 
\newcommand{\creedPath}{../Ordinaries/credo/1/credo}
% 
%%% Local Variables:
%%% mode: latex
%%% TeX-master: "missalette"
%%% End:


\begin{document}
\selectlanguage{british}

\titlePage{Haggadah} 
%\includegraphics{cover.png}
%
\section{Contents}
	\begin{itemize}
		\color{sectioncolour}
		\item Neirot Candle-lighting
		\item Kaddeish Kiddush cup 1
		\item Urchatz Handwashing (no blessing)
		\item Karpas Green herbs (dip parsley in saltwater)
		\item Yachatz Breaking the middle matzah and hiding the afikomen
		\item Maggid: beginning/intro
		\item Lachma blessing Bread of Afflictiion and invitation to the seder
		\item Ma Nishtana/Arba’ah 4 questions/4 children
		\item Haggadah Story, 10 plagues, cup 2, dayyeinu, r.gamliel’s 3 things
		\item Rachtzah handwashing (blessed)
		\item Pesach Passover Sacrifice?
		\item Matzah Unleavened Bread blessing and eat
		\item Maror horseradish/Bitter herbs
		\item Koreich Hillel Sandwich
		\item Shulchan Oreich Meal
		\item Tzafun Afikomen hunt
		\item Bareich Grace after meals
		\item Hallel Praise, door for elijah, cup of elijah and miriam’s cup
		\item Nirtzah Conclusion (next year in jerusalem)
	\end{itemize}
	%
\noindent
{\color{cyan} \rule{\linewidth}{0.5mm}}
% 

\section{Neirot}
\red{We normally begin after Havdalah by lighting the festival candles, but today we will come back to this once the sun has set fully, so we can sleep on time.}

\begin{multicols}{2}
	\CandlelightingBlessingTransliterated
	
	\columnbreak 
	\CandlelightingBlessingEnglish
	\end{multicols}

% 

\section{Kaddeish}
{I need to add ann enter after this line}

\begin{multicols}{2}
	\WineBlessingEnglish
	
	\columnbreak \WineBlessingTransliterated
	\end{multicols}
\ShehechyanuBlessing
  
\section{Urchatz}
\textit{We wash our hands without a blessing}

% 
\section{Karpas}
\textit{We dip Parsley into salt water to symbolise the bitterness, sweat and tears of the slaves}
 
\begin{multicols}{2}
	\KarpasBlessingEnglish
	
	\columnbreak \KarpasBlessingTransliterated
\end{multicols}
% 


\section{Yachatz}
\begin{multicols}{2}
	\HaMotziBlessingTransliterated
	
	\columnbreak \HaMotziBlessingEnglish
\end{multicols}

\textit{the following passage is in Aramaic:}

\begin{multicols}{2}
	\HaLachmaTransliterated

	\columnbreak \HaLachmaEnglish
\end{multicols}
% 
\section{Maggid}
\red{The creed is sung:}

\section{Ma Nishtana/Arba'ah}
\red{Immediately after the Celebrant chants the previous “Oremus”:}

\oneColumnTranslation{\offertoryTranslation}

\red{A motet or hymn may be sung here.}

\section{Haggadah}
\black{\secret}

\section{Rachtzah}

\section{Pesach}
\red{Immediately after the Celebrant concludes the chanting of the Preface:}

\section{Matzah}

\section{Maror}
\red{The Celebrant alone chants the “Pater Noster.” The Choir sings the conclusion:}
\gregorioscore{sed-libera-nos}

\section{Koreich}


\section{Shulchan Oreich}
\red{After the proceeding ``Et cum spiritu tuo:''}

\section{Tzafun}
\red{Once the Celebrant begins to distribute Holy Communion:}

\oneColumnTranslation{\communionTranslation}

\red{A motet or hymn may be sung here.}

\section{Bareich}
\black{
  \oremus
  \postcommunion
}


\section{Hallel}

% 
\section{Nirtzah}
\red{At the conclusion of Mass the Marian antiphon proper to the season is sung:}
% 

\red{A motet or hymn may follow.}
\begin{center}
  \+
\end{center}


\vfill
Typeset by Sarah Morris in ten-point Palatino.  Translations from the
MyJewishLearning.com SOURCES, and my own.

% 
The latest version of this booklet can always be found at
\url{ https://github.com/ }. Comments?
Suggestion? Found a mistake? Open an issue in the repository.
% 

{\centering\footnotesize\texttt{Compositum \today\ hora \currenttime}\par}

\end{document}

%%% Local Variables:
%%% mode: latex
%%% TeX-engine: luatex
%%% End:
